\section{Decision core}\label{sec:decision}
Part~1 defines who is acting, on what, how a decision is made, how it is bound to an operation, and how the authority state is released, stored, audited and kept available. It does not depend on storage, documents or encryption.

\subsection{Identity and trusted context}
\req{ID-001}{R1-A}{T-03}{Identity comes only from verified OIDC or OAuth tokens, workload certificates, signed S3 credentials, or another explicitly registered authentication provider. Issuer, audience, signature, expiry and intended use are validated.}
\req{ID-002}{R1-A}{T-03}{The platform \mustnot\ accept an agent- or caller-supplied role, tenant, delegation, device state or user identifier as authority.}
\req{ID-003}{R1-A}{T-03}{Adapters submit references and operation facts. The context resolver loads roles, ownership, labels, grants, session limits and approved external signals from authoritative state.}
\req{ID-004}{R1-A}{T-03}{Every trusted context field records its source, observation time, maximum age and revision where available. A missing or stale mandatory envelope fact (verified identity, tenant binding, delegation, version binding) rejects the operation.}
\req{ID-005}{R1-A}{T-03}{External lookups complete before evaluation under a bounded deadline. Policy expressions \mustnot\ fetch data or invoke tools.}
\req{ID-006}{R1-A}{T-03}{Managed-device and recent-authentication facts are usable by policy only when a trusted provider supplies them. A declared purpose may be recorded and constrained but \mustnot\ be inferred as intent. Network origin is trusted only through configured proxies that strip and regenerate forwarding headers.}
\req{ID-007}{R1-A}{T-03}{An ID token \mustnot\ be accepted as an API access token. A JWT access token is accepted only under the resource server's explicit token profile, and opaque tokens only through the provider's supported introspection flow. One JOSE implementation serves all token and envelope signature checks.}

\subsection{Delegation}
The \term{actor} is the executing workload or user. The \term{subject} is the user or service whose authority it exercises. An agent acting for a person carries both; neither inherits the other's full rights.
\req{DLG-001}{R1-A}{T-03}{A delegation record binds actor, subject, audience or executor, resource and action scope, expiry, parent delegation and an optional budget.}
\req{DLG-002}{R1-A}{T-03}{A child delegation can only narrow its parent. Delegation depth is bounded; cycles and unknown parent revisions are rejected.}
\req{DLG-003}{R1-A}{T-01, T-03}{The business-policy principal is the represented subject. Without a delegation the authenticated actor is the principal. Actor and subject role sets \mustnot\ be merged.}
\req{DLG-004}{R1-A}{T-03}{The actor's execution capability and the delegation ceiling are validated independently of business policy. Effective authority is the intersection of actor capability, delegation scope and subject permission.}
\req{DLG-005}{R1-A}{T-03, T-13}{Revocation advances a delegation generation that is checked at admission and at commit.}
\req{DLG-006}{R1-A}{T-03}{A static API key shared by several agents does not identify the user behind each call. A connector using one \mustnot\ claim user-specific enforcement without a verified delegation channel.}
\req{DLG-007}{R1-A}{T-03}{Default credentials carry no implicit administrator or cross-tenant grant. Cross-tenant access requires an explicit, reviewed sharing construct and \mustnot\ result from a missing tenant filter.}

\subsection{Resources, actions and canonical identity}
A \term{check} is one action on one resource binding. A compound operation is several checks evaluated in one batch, all of which must succeed (\rr{RES-006}).
\begin{longtable}{@{}P{0.25\textwidth}P{0.25\textwidth}P{0.31\textwidth}P{0.08\textwidth}@{}}
\toprule Resource type & Example actions & Protected facts & Phase\\\midrule\endhead
\code{agent.tool} & discover, invoke, configure & Registry identity, schema digest, connector, capability class & R1-B\\
\code{network.destination} & connect, send & Approved hostname or service identity, port, destination class, clearance & R1-B\\
\code{business.record} & read, update, approve, export & Organization, owner, workflow state, sensitive fields & R1-B\\
\code{document.version} & read-units, raw-read, discover & Canonical document and version, immutable manifest, full-content coverage & R1-C\\
\code{document.text_unit} & read, use-in-model, export-key & Canonical unit and object version, label, lineage & R1-C\\
\code{document.derived} & deliver, declassify, relabel & Lineage, effective label, intended recipient & R1-C\\
\code{text.entry} & read, use-in-model, update & Versioned text unit, parent collection & R1-C\\
\code{storage.object} & read, stat, discover, new-version, delete & Tenant, owner, project, label, immutable version, publication state & R1-C\\
\code{storage.bucket} & create-object, list, manage & Retention and encryption requirements, namespace restrictions & R1-C\\
\code{ai.model} & submit, configure & Approved provider and model, endpoint, credential binding & R1-C\\
\code{platform.policy} & propose, test, approve, activate & Bundle digest, reviewer, tenant scope, activation generation & R1-A\\
\code{crypto.recipient} & enroll, activate, receive-key, revoke & Principal or workload binding, immutable key fingerprint, purpose, status & R2\\
\code{crypto.envelope} & create, release, rewrap, retire & Object, tenant and version binding, wrapping-key reference, status & R2\\
\code{storage.replica} & repair, inspect, retire & Blob identity, node incarnation, placement generation & R3\\
\bottomrule
\end{longtable}
\req{RES-001}{R1-A}{T-01, T-06}{Every protected operation maps to one or more (resource, action) checks from the registry. An operation with no registered mapping fails closed.}
\req{RES-002}{R1-A}{T-06}{A canonical resource registry maps S3 bucket, key and version, MCP URI and tool selectors to the same tenant-scoped resource UUID, immutable content version, policy kind, version and scope, and trusted attributes. Role membership and principal context resolve through one identity pipeline.}
\req{RES-003}{R1-A}{T-01, T-06}{Content reading uses the canonical action \code{read}. Legacy route names such as \code{storage.object.read} are aliases, not competing policies. A caller \mustnot\ obtain a less restrictive policy kind or scope by changing protocol.}
\req{RES-004}{R1-A}{T-06}{Alias resolution happens once under a declared catalog snapshot and is bound to the operation. The resource generation is revalidated before admission and publication. ``Latest'' \mustnot\ be resolved again after authorization.}
\req{RES-005}{R1-A}{T-06}{A rename changes a namespace alias, not the resource's permission identity. An edit creates a new immutable version.}
\req{RES-006}{R1-A}{T-06}{A compound operation requires a separate allowed check for every resource and action it actually needs. One allow \mustnot\ authorize an unrelated effect, and \code{tool.invoke} alone \mustnot\ authorize reading the tool's source object or sending to its destination.}
\req{RES-007}{R1-A}{T-01, T-06}{A channel difference affects a decision only through an explicitly declared, trusted context attribute that the policy deliberately uses. Authentication or unsupported-protocol failures stay separate restrictions and \mustnot\ turn a business denial into an allow.}
\req{RES-008}{R1-A}{T-06}{Parent-resource checks configured for a content family apply identically through every transport as explicit compound checks, not as a hidden change to scope inheritance.}

\subsection{Policy evaluation}\label{sec:policy}
\subsubsection{Evaluator contract}
The decision service calls the business evaluator through a narrow contract so that the engine behind it (decision D-01) can be native or upstream.
\begin{lstlisting}
// Logical contract. Pure: no I/O, no ambient state, time passed in.
EvaluatorRequest {
  snapshot_digest,
  policy_set: BUSINESS | CEILING,
  principal: {id, roles, attrs},
  resource:  {kind, id, version, scope, attrs},
  actions:   [string],
  aux:       {envelope facts mapped as documented attributes},
  now
}
EvaluatorResult {
  per_action: [{action, effect, outputs, metadata, diagnostics}],
  errors:     [{action, kind}]          // strict mode: affected action is DENY
}
\end{lstlisting}
\req{POL-001}{R1-A}{T-01}{Business policy evaluation sits behind the evaluator contract. The conformance corpus (\rr{TEST-003}) applies to any engine that implements the contract.}
\req{POL-002}{R1-A}{T-01, T-02}{The evaluator implements the project's own policy format and evaluation rules of Section~\ref{sec:policyformat}. A document or construct that the format does not define is rejected at bundle validation and \mustnot\ be silently ignored.}
\req{POL-003}{R1-A}{T-01}{Evaluation proceeds in this order: (1) authenticate and validate the execution envelope; (1b) evaluate the ceiling set for every check, as an always-run gate that does not depend on the business result; (2) normalize resource, action and principal input and resolve the scope chain; (3) evaluate the business set by the rules of Section~\ref{sec:policyformat}, in which a matching deny always overrides an allow; (4) apply the business result unchanged to that check and combine it with the ceiling result so that a ceiling \code{DENY} always wins; (5) fulfil platform prerequisites and obligations. Platform steps add rejections only and \mustnot\ turn a business denial into an allow.}
\req{POL-004}{R1-A}{T-01}{The response records the business decision and any platform rejection reason separately. A protocol or authentication failure \mustnot\ be reported as a policy-combination rule, and no adapter implements alternate granting logic.}
\req{POL-005}{R1-A}{T-02}{Evaluation is strict. An error in the condition of any candidate rule for a requested action yields \code{DENY} for that action, or an objection in the ceiling set. The format has no lenient mode, so no configuration can select one.}
\req{POL-006}{R1-A}{T-02}{A missing business attribute is handled by the production profile and \mustnot\ fall through to an allow. Trusted-envelope facts never depend on business-schema settings: invalid credentials, wrong tenant or resource binding, unavailable authority, malformed transport, process faults and exhausted limits always reject.}
\req{POL-007}{R1-A}{T-01, T-02}{Platform ceilings are a separate, deny-only ceiling set in the signed bundle. Every check is also evaluated against the ceiling set through the evaluator contract (\code{policy_set: CEILING}) before and independently of the business evaluation, and a ceiling \code{DENY} or a ceiling evaluation error denies the check whatever any principal, role, derived-role or scope result says. Bundle validation rejects a ceiling set that contains a principal policy, an allow rule or a scope override.}
\req{POL-008}{R1-A}{T-01}{Time is captured freshly for each request or batch and never frozen at bundle preparation. Admission and commit checks use current server time.}
\req{POL-009}{R1-A}{T-02}{Policy expressions cannot fetch network data, read files, execute scripts or acquire credentials. Compilation, input size, evaluation cost, memory and concurrency are bounded independently of business-policy options.}
\req{POL-010}{R1-A}{T-01}{A snapshot is prepared once per immutable bundle. It contains the validated policy graph, schemas, checked executable expressions, configuration and dispatch indexes keyed by static resource, action, version and scope facts. The index \mustnot\ omit a rule because another path looks permissive. A compiled expression is shared only when the complete environment identity (types, function signatures, schema digests, library and options versions, cost model, evaluator-semantics version) matches.}
\req{POL-011}{R1-A}{T-01}{Policy outputs are data. A mapping from selected output fields to platform obligations is allowed only when explicitly configured, versioned and validated. No mandatory proprietary policy field is introduced for approvals or storage behavior.}
\reqn{POL-012}{R1-A}{T-01}{The conformance corpus of policy-evaluation vectors is part of this specification and runs against the engine in CI. No upstream engine is used as a reference or oracle (D-01).}
\req{POL-013}{R1-A}{T-02}{YAML authoring is converted offline into the same strict typed representation as JSON. Ambiguous or unsafe constructs are rejected, and evaluation never runs a YAML parser.}
\reqn{POL-014}{R1-A}{T-01}{Scopes form a dot-separated chain from the request's scope to the root. Rules from every scope on the chain are combined by the same deny-overrides rule, so a child scope can add allows but cannot lift a deny from a parent scope or from the root.}

\subsubsection{Policy format}\label{sec:policyformat}
The policy format is the project's own (D-01). A document holds rules for one resource type in one scope and one policy set:
\begin{lstlisting}
{
  "format": "fga.policy/v1",
  "set": "business",
  "resource": "document.text_unit",
  "scope": "acme",
  "rules": [
    {"id": "reviewers-read-own-projects",
     "effect": "allow",
     "actions": ["read"],
     "roles": ["reviewer"],
     "when": "resource.attributes.project_id in principal.attributes.project_ids"}
  ]
}
\end{lstlisting}
The ceiling set uses the same document shape with \code{"set": "ceiling"} and deny rules only, for example a rule that denies \code{read} for any role \code{when} the unit is not LIVE.

For one principal, resource and action, the engine:
\begin{enumerate}
\item collects the candidate rules of the requested set for the resource type from the request's scope and each ancestor up to the root;
\item treats a rule as matching when it lists the action (or \code{*}), the principal holds one of its roles (or \code{*}), and its condition is true or absent;
\item for the business set, returns \code{DENY} if any matching rule denies, otherwise \code{ALLOW} if any matching rule allows, otherwise \code{DENY};
\item for the ceiling set, reports an objection if any matching rule denies, and no objection otherwise;
\item denies (or objects) when the condition of any candidate rule fails to evaluate.
\end{enumerate}
An allow rule in the business set may list prerequisites in \code{requires} (\code{approval}, \code{reservation}). An allowed check carries the union of the prerequisites of every matching allow rule, and the decision is \code{NEEDS_APPROVAL} or \code{NEEDS_RESERVATION} until the prerequisite is satisfied through the control plane (\rr{APR-001}, \rr{APR-007}).

Conditions are CEL expressions over \code{principal} (\code{id}, \code{roles}, \code{attributes}), \code{resource} (\code{type}, \code{id}, \code{scope}, \code{attributes}) and \code{request} (\code{now}). They are compiled once per bundle and run under a cost limit. Because a deny always overrides an allow, holding an extra role cannot undo a deny, and no business rule can bypass a ceiling. Principal-specific policies, derived roles, shared variables and outputs are not part of version 1; they can be added later without changing the meaning of existing policies.

\subsection{Decision API}\label{sec:decisionapi}
\begin{lstlisting}
DecideBatch request {
  request_id, tenant_id, adapter_id, operation_id,
  actor:      {id, type, identity_revision},
  subject:    {id, type, identity_revision},
  delegation: {id, revision},
  session:    {id, binding_revision, context_key,
               mark_revision, label_order_generation},
  checks:     [{check_id, action,
                resource: {type, id, version, generation, metadata_revision}}],
  operation_digest, required_policy_epoch, deadline
}
DecisionResponse {
  outcome: ALLOW | DENY | NEEDS_APPROVAL | NEEDS_RESERVATION | ERROR,
  reason_code, decision_id,
  policy_results: [{check_id, principal_id, resource_id, action,
                    business_effect | NOT_EVALUATED, ceiling_effect,
                    policy_outputs,
                    policy_metadata, diagnostics}],
  bundle_digest, policy_epoch, evaluator_version,
  actor_id, subject_id, delegation_revision,
  operation_digest, resource_bindings_digest, context_revisions,
  evaluated_at, valid_until,
  obligations: [{type, version, parameters_or_ref, parameters_digest}],
  session_mark_revision_after,
  signature                       // detached JWS, see DEC-008
}
\end{lstlisting}
\code{Decide} is a batch of one. A platform-gate rejection returns \code{DENY} with a platform reason code and \code{NOT_EVALUATED} business effects. \code{ERROR} means the service could not decide (timeout, unavailable authority, resource limit).
\req{DEC-001}{R1-A}{T-01}{The decision API is exposed as versioned gRPC for internal services and an authenticated JSON facade for SDKs. The caller \mustnot\ choose an action mapping different from the registered operation.}
\req{DEC-002}{R1-A}{T-01}{The enforcement point acts only on an explicit, complete \code{ALLOW} bound to its exact operation. \code{DENY}, \code{NEEDS_APPROVAL}, \code{NEEDS_RESERVATION}, \code{ERROR}, transport errors, malformed data, unknown required obligations and missing results all stop execution. A prerequisite outcome is satisfied through the control plane, after which a new \code{ALLOW} whose context revisions include the approval or reservation is required.}
\req{DEC-003}{R1-A}{T-01}{Each response carries the per-check business and ceiling effects, outputs, metadata, redacted diagnostics, bundle digest, policy epoch, evaluator version, identity and delegation revisions, operation and resource-bindings digests, validity and obligations. A platform gate that stops evaluation is reported with an explicit not-evaluated marker and \mustnot\ authorize execution.}
\req{DEC-004}{R1-A}{T-01}{A batch captures one policy snapshot and one authority envelope. Items have unique identifiers and exact resource, action and revision bindings. A missing, duplicated or misbound result \mustnot\ authorize its operation, and a failed item \mustnot\ grant its peers. A batch is an evaluation optimization, not a transaction over writes or external effects.}
\req{DEC-005}{R1-A}{T-13}{\code{valid_until} is bounded by the earliest of identity or session expiry, delegation expiry, approval or grant expiry, context freshness limit, request deadline and a conservative decision-age limit. A decision is not a general bearer permission.}
\req{DEC-006}{R1-A}{T-13}{At read admission and immediately before a mutation's irreversible or egress checkpoint, the executor rechecks expiry and every authority-bearing revision recorded with the decision: resource generation; metadata revision; access and label revisions of the parent and of each selected unit; delegation and identity revisions; context mark revision and label-order generation; and the required policy epoch. A mismatch requires a new decision or rejection. The executor uses the same canonical operation that was evaluated.}
\req{DEC-007}{R1-A}{T-03}{Identity and capability checks are enforced at the decision API itself. mTLS workload identity or an equivalent bound credential identifies the adapter, and an ordinary user cannot submit fabricated verified context.}
\reqn{DEC-008}{R1-A}{T-03}{The decision service signs each \code{ALLOW} record with a registered decision-signing key, used for no other purpose, whenever the decision will be used across a process or role boundary or retained. The signed bytes are one canonical serialization (strict JSON with no unknown or duplicate fields, signed exactly as transmitted) containing: decision ID; audience (deployment, role and service); tenant; actor, subject, delegation and context identity with their revisions; operation digest; every check with its resource binding and effect; bundle digest and policy epoch; access, label and context-mark revisions and the label-order generation; and issue time and \code{valid_until}. A verifier checks signature, audience, validity and the replay rule, and still performs the fresh revision checks of \rr{DEC-006}: a signature does not replace them. A role \mustnot\ accept a decision whose audience is another role. A signed decision proves that the sole authority evaluated the request it describes; it is not proof that the user's intent or the request's context was validly established.}
\req{DEC-009}{R1-A}{T-22}{Initial qualification limits are a 250 ms single-decision budget, 64 KiB of context, 50 resources per batch, 8 KiB of expression text, bounded AST depth, 10,000 expression cost units per condition and a total request budget. These are proposed limits, not measured guarantees; independent input, compilation, allocation, concurrency and cancellation limits are mandatory, and cancellation reaches both queued and running evaluation.}
\req{DEC-010}{R1-A}{T-01}{Error diagnostics distinguish a policy denial from a service failure without exposing protected details to unauthorized clients.}
\reqn{DEC-011}{R1-A}{T-03}{Consumption of a single-use signed decision (mutation and key release) is an atomic, durably acknowledged compare-and-set in the catalog, shared by every replica of the verifying role and recorded before the effect begins. A replica that cannot reach the consumption store refuses the decision.}
\reqn{DEC-012}{R1-A}{T-03}{Decision-signing keys are distinct from workload, issuer and encryption keys. Key identifiers reach verifiers through the control plane, rotation uses an overlap window no shorter than the longest \code{valid_until} plus clock skew, and a compromised signing key is revoked by a revision that verifiers must observe before accepting further decisions. Rotation and revocation are qualified under T-03.}

\subsection{Policy releases and readiness}
\req{REL-001}{R1-A}{T-14}{A policy release is an immutable, signed bundle whose manifest includes schema and evaluator versions, source hashes, the business set, the ceiling set, the label order and its generation, resource schemas, operation registry, pure-function allowlist, expression-library version and cost model, test evidence, target audience and digest. Each instance verifies, parses, compiles, validates and checks a candidate into an immutable in-memory snapshot before atomically replacing its active pointer. A failed activation preserves the prior valid state and does not acknowledge the candidate.}
\req{REL-002}{R1-A}{T-14}{Activation records the states RECEIVED, VERIFIED, PREPARED, READY, ACTIVE or REJECTED with digest, environment version, preparation error and timestamps. Every decision uses exactly one snapshot, and preparation and old-snapshot coexistence are resource-bounded.}
\req{REL-003}{R1-A}{T-13, T-14}{A node that cannot establish the authoritative required epoch \mustnot\ authorize new operations. The last valid snapshot is kept only while it satisfies that epoch. Offline authorization is not an availability fallback.}
\req{REL-004}{R1-A}{T-14}{The control plane stages a candidate pool, verifies sentinel allow and deny cases, records ready instances by digest, advances the durable tenant-required epoch, and routes only to compatible instances. A restrictive activation may stop admissions and drain bounded mutation leases. A rollback is a new reviewed release with a higher epoch, because silently reverting may restore revoked access.}
\reqn{REL-005}{R1-A}{T-14}{Before activation the candidate is run over the conformance corpus and the tenant's sentinel cases, and the change in effect against the active release is reported to the reviewer, who must acknowledge it. The report covers the corpus and sentinels only; it does not enumerate every principal.}
\req{REL-006}{R1-A}{T-14, T-19}{Liveness reports process progress and does not restart a healthy process because PostgreSQL, KMS or a peer is unavailable. Readiness is role- and capability-specific (policy digest and epoch, catalog authority, workload credential validity, required crypto and connector dependencies). Probes return redacted aggregate state, are rate-limited, and do not create a KMS key. A readiness result is never a grant; requests perform their own admission checks.}

\subsection{State and catalog}
PostgreSQL holds authoritative identities, grants, resource bindings, revisions, approvals, reservations and audit-outbox records.
\begin{longtable}{@{}P{0.20\textwidth}P{0.62\textwidth}P{0.08\textwidth}@{}}
\toprule Entity & Minimum authoritative fields & Phase\\\midrule\endhead
Principal, delegation & Stable IDs; tenant; actor and subject bindings; credential class; roles; limits; expiry; revocation and identity revisions & R1-A\\
Resource & UUID; type; tenant; protected attributes; metadata revision; resource generation; lifecycle state & R1-A\\
Policy release & Canonical manifest digest; schema, configuration, evaluator and expression-library versions; tests; publisher signature; activation epoch & R1-A\\
Audit outbox & Unique event ID; decision and operation IDs; redacted bindings; outcome; delivery status & R1-A\\
Tool registry & Stable tool ID; schema digest; connector version; operation mapper; allowed destinations; required capabilities; credential binding & R1-B\\
Invocation & Tenant and idempotency unique key; argument digest; state; dispatch attempt and fencing token; downstream receipt; reconciliation status & R1-B\\
Approval & Exact operation digest; approver; subject and actor; resource revisions; expiry; allowed uses; consumption state & R1-B\\
Reservation & Budget or grant ID; minor-unit amount or use count; invocation ID; status; lease; reconciliation state & R1-B\\
Object, version & Bucket and exact key; stable object UUID; current committed version; content generation; immutable blob ID; length and checksums & R1-C\\
Document manifest, text unit & Document and version UUIDs; ordered unit IDs; manifest digest; extraction and normalization version; unit order; labels; access revision; provenance & R1-C\\
Upload intent & Operation and idempotency IDs; principal and delegation; expected generation and current version; intended labels; staged blob; state; expiry & R1-C\\
Context mark & Context key (runtime context ID, else delegation ID); bound actor, subject and delegation; label-order generation; highest effective label; mark revision & R1-C\\
Recipient key, key envelope & Tenant; principal or workload; immutable key ID and fingerprint; purpose and suite; activation and revocation revision; wrapped keyset; generation; status & R2\\
Maintenance job & Kind; tenant and blob scope; expected placement generation; idempotency key; lease and fence; progress; initiating actor & R3\\
\bottomrule
\end{longtable}
\req{DAT-001}{R1-A}{T-01}{All authoritative lookups include tenant scope. Tenant-consistent references are enforced with database constraints or equivalent transaction checks. Tenant isolation \mustnot\ depend on policy authors writing a tenant condition.}
\req{DAT-002}{R1-A}{T-03}{Protected fields change only through explicit authorized control operations with expected revisions. Raw S3 tags, tool arguments and external payloads \mustnot\ become trusted ownership or classification records.}
\req{DAT-003}{R1-A}{T-06, T-13}{Resource metadata revision and content or resource generation are distinct. Every logical publication, deletion or target-state mutation increments the generation, and approval and execution comparisons include state changes even when access labels have not changed.}
\req{DAT-004}{R1-A}{T-13, T-14}{Permission-critical facts use authoritative consistent reads. The first release has no allow-result cache. A cache, if added, requires a maximum age, revision-aware keys and explicit invalidation. A decision uses one immutable policy snapshot and one declared context snapshot. Search indexes and storage notifications are discovery aids, not access authority.}
\req{DAT-005}{R1-B}{T-05}{Money uses integer minor units with an explicit currency. Binary floating point \mustnot\ be used for permission thresholds or budget consumption.}

\subsection{Audit}
\req{AUD-001}{R1-A}{T-15}{Decisions and actual outcomes are logged separately and joined by decision, operation and invocation IDs. Records include actor, subject, tenant, action, resource, version and generation, relevant context provenance, policy digest and epoch, reason, latency and redacted obligations. External tools record dispatch and receipt or reconciliation evidence.}
\req{AUD-002}{R1-A}{T-15}{Attempt, decision, key release, durable mutation, approval consumption and external delivery are distinct events. Mandatory records are appended transactionally to an outbox and \mustnot\ silently disappear when the exporter is disabled or a queue is full; capacity alarms and declared admission limits act before required evidence can no longer be recorded. Operational logs and optional notifications are not a substitute for the outbox.}
\req{AUD-003}{R1-A}{T-19}{Audit delivery is idempotent with explicit retention and access policies. Sensitive labels and arguments are minimized or protected, lag, loss, stale policy nodes, failed context resolution and revocation latency are monitored, and metrics \mustnot\ create a cross-tenant information channel.}
\req{AUD-004}{R1-A}{REV}{A key-release audit record proves that a key was issued, not that a user opened plaintext or complied with a purpose restriction. Documentation \mustnot\ present it as proof of use.}

\subsection{Administration}
\req{ADM-001}{R1-A}{T-20}{Policy authors, reviewers, activators, key administrators, storage operators and data users are separate roles. A policy administrator is not automatically a data reader, and sensitive policy or key changes use the approval and reservation model where configured.}
\req{ADM-002}{R1-A}{T-20}{An ordinary API request \mustnot\ disable the authorizer, change trusted issuers, relax workload audience checks, replace a key provider or grant itself policy-administrator authority. Those changes use explicit administrative actions and reviewed control transactions.}
\req{ADM-003}{R1-A}{REV}{Bootstrap and recovery authority sits outside ordinary agent paths, with controlled local or separately authenticated recovery and a durable incident audit. Recovery procedures \mustnot\ claim to defend against a host operator who can replace the running software.}

\subsection{Availability of authority state}\label{sec:availability}
This section covers the decision, control and catalog services. Blob replication is in Part~5.
\begin{longtable}{@{}P{0.20\textwidth}P{0.40\textwidth}P{0.30\textwidth}@{}}
\toprule Layer & Initial production topology & Failure contract\\\midrule\endhead
Gateways, decision, tools & At least 3 stateless replicas across 3 failure domains, required bundle digest on each & Load balancer removes unready replicas; stale or unavailable authority never becomes allow\\
Tool executors & At least 3 workers sharing durable invocation state & Leases and fencing plus downstream idempotency; no blind redispatch after timeout\\
Catalog PostgreSQL & Primary plus 2 standbys; commit durable on the primary and at least one synchronized remote standby (quorum commit) & HA coordinator promotes only a safe candidate and fences the old primary; no silent asynchronous downgrade\\
\bottomrule
\end{longtable}
\req{AVL-001}{R1-C}{T-15, T-16}{Production runs the topology above. Single-node operation is a development profile and is rejected by production configuration.}
\req{AVL-002}{R1-C}{T-16}{Authorization-critical reads use the primary or a standby proven sufficiently fresh by a log-position check. An arbitrarily lagging standby \mustnot\ restore revoked grants.}
\req{AVL-003}{R1-C}{T-16}{PostgreSQL runs with \code{fsync=on}, \code{full_page_writes=on}, \code{synchronous_commit=on} or stronger, and a non-empty synchronous standby set; \code{remote_write} alone is insufficient against an operating-system crash. With Patroni, strict synchronous mode is enabled. Failover uses a managed HA service or established controller with quorum coordination and old-primary fencing, and promoting an unqualified stale replica voids the zero-loss objective.}
\req{AVL-004}{R1-C}{T-16}{A timeout or cancellation while awaiting a synchronous commit is an unknown outcome, and a lookup by operation ID on the primary is not evidence of durability: after a cancelled wait the transaction can be visible on the primary without having been replicated. Before returning durable success, dispatching an external effect or releasing a key, the platform obtains proof that the required synchronous standbys have flushed WAL at or beyond the transaction's commit position on the current fenced timeline, either from an acknowledged synchronous commit ordered after it or from the replication flush position. Until then the outcome stays unknown.}
\req{AVL-005}{R1-C}{T-15, T-18}{The design target is zero loss of acknowledged operations under the loss of one failure domain in a crash-fault model, conditional on synchronous database acknowledgement, correct fencing and available key state; it is not a guarantee against simultaneous multi-domain loss, Byzantine replicas, malicious deletion or untested hardware. Recovery of the target service within five minutes is an acceptance objective, not a benchmark. Encrypted, access-separated backups of the catalog, policy bundles and (from R2) key envelopes and KMS recovery state are restored together in tests, and replication is not a substitute for backups.}
\req{AVL-006}{R1-C}{T-22}{Tenant-aware admission limits cover requests, bytes, active streams, expression cost, scan candidates, reservations, tool concurrency and KMS operations. Overload uses backpressure and bounded queues rather than unbounded goroutines, and retries use jitter and budgets. Circuit breakers reject rather than grant.}
\req{AVL-007}{R1-C}{T-22}{The first metadata authority is one PostgreSQL primary with HA. When measured write capacity is exceeded, tenants are sharded across independent catalog groups with an explicit tenant-to-shard registry and migration fencing, and cross-shard grants and budgets are not treated as local atomic operations. The first region is active for each tenant; multi-region active-active authorization requires a later consistency design.}
